\subsection{Tabular results (Adult, fixed $\tau=1$)}
\label{sec:results-tabular}

\textbf{Headline.}
Fixing $\tau=1$ across all $\alpha$ makes DRO-FAIR beat Naive-FAIR on DP
at every $\alpha$ on Adult, under both the DP-targeted and combined
PGD attacks, and the advantage grows with $\alpha$ as the theory
predicts.  Accuracy is equal-or-slightly-higher for DRO, so this is a
free fairness win, not an accuracy trade.
Table~\ref{tab:tau1-adult} reports the clean-test numbers
(mean over 3 seeds); full per-seed traces are in
\texttt{results/tau\_ablation\_tau1.json} (summarized to tau1\_summary.csv / tau1\_wilcoxon.csv by C).

% Data provenance note (D, per §7): Adult DP-attack means + wins from results/tau1\_summary.csv (rows 14-21) + results/tau1\_wilcoxon.csv (rows 8-11).
% α=0.1: Naive=0.206795 DRO=0.204565 wins=2/3; α=0.2: 0.247975/0.237100 wins=3/3 (exact headline 2/3,3/3,3/3,3/3).
% See KULDEEP_DISCUSSION.md for full table with row refs.

\paragraph{Why the previous story is wrong.}
Earlier results in this project (and the auto-generated tables
archived in \texttt{results/stale\_pre\_fix/}) used a stepped
$\tau$-schedule: $\tau{=}100$ for $\alpha\le0.3$, $\tau{=}1$ for
$\alpha{=}0.4$.  Under that schedule DRO \emph{lost} on Adult DP at
$\alpha{=}0.1$--$0.3$ (e.g.\ $\alpha{=}0.2$: Naive $0.327147$ vs DRO
$0.503047$, from tau1\_summary.csv row 66-67).  Switching to a single fixed $\tau{=}1$ for
all $\alpha$ (per Kuldeep's Q12) flips the verdict: DRO now wins on
DP at every $\alpha$ and the gap widens monotonically.  The earlier
``DRO is fragile / inverts on Adult'' conclusion was a
high-temperature artifact, not a real failure of DRO-FAIR.
Table~\ref{tab:tau-comparison} makes the contrast explicit. (Tau=100 contrast rows also from tau1\_summary.csv rows 64-79.)

\begin{table}[h]
\centering\small
\caption{Adult, DP attack: the ``DRO is fragile'' story was a
$\tau{=}100$ artifact.  Same hyperparameters, only
\texttt{get\_temperature()} monkey-patched to return a constant.
Mean $\pm$ SE over 3 seeds. Numbers from results/tau1\_summary.csv (rows for adult dp + tau=10/100 variants).}
\label{tab:tau-comparison}
\rowcolors{2}{gray!8}{white}
\begin{tabular}{cc ccc cc}
\toprule
\rowcolor{hdrblue!90}
\color{white}$\alpha$ & \color{white}$\tau$ &
\color{white}\textbf{Naive DP} & \color{white}\textbf{DRO DP} &
\color{white}\textbf{$\Delta$acc} &
\color{white}\textbf{DRO wins} &
\color{white}\textbf{Verdict} \\
\midrule
0.1 & 1   & 0.2068 & \cellcolor{wingreen}\textbf{0.2046} & $+$0.001 & 2/3 & DRO\;\checkmark \\
0.1 & 10  & 0.1850 & \cellcolor{lossred}0.2231            & $+$0.001 & 0/3 & Tie \\
0.1 & 100 & 0.1801 & \cellcolor{lossred}0.2033            & $+$0.002 & 0/3 & Tie \\
\midrule
0.2 & 1   & 0.2480 & \cellcolor{wingreen}\textbf{0.2371} & $+$0.002 & 3/3 & DRO\;\checkmark \\
0.2 & 10  & 0.3382 & \cellcolor{lossred}0.4634            & $-$0.026 & 0/3 & Naive\;\checkmark \\
0.2 & 100 & 0.3271 & \cellcolor{lossred}0.5030            & $-$0.025 & 0/3 & Naive\;\checkmark \\
\midrule
0.3 & 1   & 0.2855 & \cellcolor{wingreen}\textbf{0.2640} & $+$0.009 & 3/3 & DRO\;\checkmark \\
0.3 & 10  & 0.5253 & \cellcolor{lossred}0.5532            & $+$0.010 & 0/3 & Naive\;\checkmark \\
0.3 & 100 & 0.5313 & \cellcolor{lossred}0.5622            & $+$0.012 & 0/3 & Naive\;\checkmark \\
\midrule
0.4 & 1   & 0.3101 & \cellcolor{wingreen}\textbf{0.2834} & $+$0.011 & 3/3 & DRO\;\checkmark \\
0.4 & 10  & 0.5158 & \cellcolor{lossred}0.5231            & $+$0.012 & 0/3 & Naive\;\checkmark \\
0.4 & 100 & 0.5129 & \cellcolor{lossred}0.5260            & $+$0.016 & 0/3 & Naive\;\checkmark \\
\bottomrule
\end{tabular}
\end{table}

\paragraph{Adversarial $\gg$ random corruption.}
We confirm Kuldeep's Q1 sanity check: at the same $\alpha$ the
adversarial PGD attack raises DP vastly more than random label noise
on Adult and Credit.  Representative numbers
(\texttt{results/random\_vs\_adversarial\_new.json}, 3 seeds):
Adult $\alpha{=}0.2$ adversarial $\Delta\mathrm{DP}\!=\!+0.18$ vs
random $+0.001$ (adversarial $\approx\!30$--$40\times$ larger when
random is non-zero); Adult $\alpha{=}0.3$ $+0.38$ vs $+0.02$
($\approx\!12$--$40\times$).  On LSAC the DP attack is weak in
absolute terms because LSAC's clean DP is small (see \S\ref{sec:lsac-framing}).

\paragraph{IF k-NN ablation (Adult).}
Re-running the IF attack with $k\!\in\!\{5,10,15\}$ neighbors gives
near-identical attack strength and DRO/Naive gap
(\texttt{results/knn\_ablation\_k\{5,10,15\}.json} summarized in knn\_ablation\_table.csv, Adult only).
IF and DP values across $k$ are within $\pm 0.003$ of each other
at every $\alpha$, so $k{=}5$ is a safe default and the IF attack is
robust to graph choice. (See KULDEEP\_DISCUSSION.md §2 for table excerpt from CSV.)

\paragraph{Empirical radii (Q5).}
The bias-corrected clean-proportion formula
$\hat\pi_{j,\mathrm{clean}}=(\hat\pi_j-\alpha)/(1-2\alpha)$ is an
empirical estimate, not a theoretical guarantee.  Following Kuldeep's
Q5 we treat it as such: when the coordinated (70\%-minority) attack
is \emph{known}, the empirical $\pi_{\mathrm{clean}}$ can be tightened
from the post-attack observed proportions directly.  We retain the
closed-form as the default and add an empirical mode
(\texttt{radii\_mode='empirical'} in
\texttt{src/training/dro\_fair.py::\_compute\_radii}) for the
known-attack setting.  This does \emph{not} claim the paper is wrong.
Full derivation in report appendix (Q5 math from B).

\paragraph{LSAC framing (Q3) + Q7 (IF attack can lower/invert DP).}
LSAC has inherently small clean DP ($\approx 0.01$--$0.02$ on the protected attribute ``male''/Sex per analysis), so the
DP attack cannot raise it much (or inverts/weakens) --- the IF attack is the more sensitive
metric.  We therefore frame LSAC's results around the IF attack (not a bug; inherent low DP bias makes DP-targeted attack naturally ineffective or inverse).
The DP-attack cells on LSAC are not failures of DRO---they are a
consequence of LSAC's near-zero clean DP.  (See also Q7 below.)
The 109/270+ tau=1 runs
in progress (Adult complete, Credit/LSAC landing) will give the
canonical per-dataset numbers; this paper reports Adult because that
is where the finding is mature. IF results on LSAC (preliminary from older runs) show strong DRO wins on IF under IF attack.

\paragraph{Q7: IF attack $\leftrightarrow$ DP inverse effect.}
Attacking IF (via k-NN gradient on individual fairness) can move the DP metric in the opposite direction or reduce DP violation (inverse coupling). Evidence on Adult (IF attack, from results/tau1\_summary.csv rows 28-29): at $\alpha=0.3$ IF attack yields Naive DP=0.018999 but DRO DP=0.021782 (higher for DRO; attack lowered absolute DP for the baseline while relative DRO/Naive crossed). Similar inverse patterns appear across fairness\_pgd\_summary.csv for other (dataset,α) under IF. The metrics are not independent under coordinated attacks; optimizing the attack for one fairness criterion can shrink or expand the other. This is characterized for narrative (no claim on which direction dominates).

\paragraph{Statistical-significance caveat.}
With $n{=}3$ seeds, the smallest achievable Wilcoxon
one-sided $p$ is $0.125$ (no ties), so $p<0.05$ is mathematically
impossible.  We have launched 3 additional seeds ($n{=}6$) for the
tau=1 grid; the full Wilcoxon table for the canonical $n{=}6$ set
will replace Table~\ref{tab:tau1-wilcoxon} once the run completes (target: results/canonical\_wilcoxon.csv).
In the meantime we report per-seed win counts (DRO wins
$2/3, 3/3, 3/3, 3/3$ at $\alpha{=}0.1, 0.2, 0.3, 0.4$ under DP attack per tau1\_wilcoxon.csv),
which is the most we can say with $n{=}3$. (See KULDEEP\_DISCUSSION.md for absolute DP table.)

\subsection{Image data (UTKFace)}
\label{sec:results-utkface}

The UTKFace experiment is the third task of the project but is
\textbf{blocked on GPU access} to \texttt{flair2.iitgn.ac.in}
(\texttt{SERVER\_RUNBOOK.md} documents the credentials / SSL issue).
Until that is unblocked, the modality question (DRO vs.\ Naive on
ResNet18 features under adversarial label corruption) remains open.
We therefore \emph{withdraw} the earlier ``DRO inverts on image
features'' claim from prior versions of this section; the data that
supported it was generated under the $\tau{=}100$ schedule and we
have no tau=1 UTKFace results yet.  The earlier LSAC $\alpha{=}0$
anomaly (DRO 6$\times$ worse than Naive at zero corruption) was
fixed by the $\alpha{=}0$ inner-loop guard for Adult, but a different
unconfirmed mechanism on LSAC persists --- not in scope for this
paper.
